\documentclass[11pt]{article}

\usepackage[margin=1in]{geometry}
\usepackage{graphicx}
\usepackage{booktabs}
\usepackage{amsmath,amssymb,amsthm}
\usepackage{xcolor}
\usepackage{mdframed}
\usepackage{hyperref}
\usepackage{caption}
\hypersetup{colorlinks=true, linkcolor=blue, citecolor=blue, urlcolor=blue}

\theoremstyle{plain}
\newtheorem{proposition}{Proposition}
\theoremstyle{remark}
\newtheorem{remark}{Remark}
\newtheorem{assumption}{Assumption}

\DeclareMathOperator{\KL}{KL}
\DeclareMathOperator{\MMD}{MMD}
\newcommand{\R}{\mathbb{R}}
\newcommand{\E}{\mathbb{E}}
\newcommand{\dd}{\,\mathrm{d}}

% --- plain-language callout box (supporting intuition only) -----------------
\definecolor{boxbg}{RGB}{237,243,252}
\definecolor{boxln}{RGB}{120,160,210}
\newmdenv[backgroundcolor=boxbg,linecolor=boxln,linewidth=0.6pt,
  skipabove=9pt,skipbelow=9pt,
  innertopmargin=7pt,innerbottommargin=7pt,
  innerleftmargin=10pt,innerrightmargin=10pt]{plainbox}
\newcommand{\plainwords}[1]{\begin{plainbox}\noindent\textit{Intuition.\ }#1\end{plainbox}}

% --- robust includes: skip gracefully if an artifact is missing -------------
\newcommand{\figbox}[3]{%
  \IfFileExists{figures/#1.pdf}{%
    \begin{figure}[t]\centering
      \includegraphics[width=\linewidth]{figures/#1.pdf}
      \caption{#2}\label{#3}
    \end{figure}}{\par\noindent\textit{[figure \texttt{\detokenize{#1}} not found]}\par}}
\newcommand{\figboxw}[4]{%
  \IfFileExists{figures/#1.pdf}{%
    \begin{figure}[t]\centering
      \includegraphics[width=#4\linewidth]{figures/#1.pdf}
      \caption{#2}\label{#3}
    \end{figure}}{\par\noindent\textit{[figure \texttt{\detokenize{#1}} not found]}\par}}
\newcommand{\inputtab}[1]{\IfFileExists{tables/#1.tex}{\input{tables/#1.tex}}{}}

% --- live headline numbers (regenerated by build_report.py) -----------------
\IfFileExists{numbers.tex}{\newcommand{\MogLSBcovK}{1.00}
\newcommand{\MogLSBemcK}{0.94}
\newcommand{\MogPTcovK}{0.23}
\newcommand{\MogPTemcK}{0.07}
\newcommand{\MogHMCcovK}{0.18}
\newcommand{\MogLSBwtwoK}{10.00}
\newcommand{\MogPTwtwoK}{33.85}
\newcommand{\MogLSBemcSmall}{0.92}
\newcommand{\MogLSBemcBig}{0.83}
\newcommand{\MogLSBruntimeBig}{184.5}
\newcommand{\MogPTruntimeBig}{2.6}
\newcommand{\MogLSBpotBig}{7.7}
\newcommand{\MwLSBcemcBig}{0.96}
\newcommand{\MwPTcemcBig}{0.00}
\newcommand{\MwLSBcemcSmall}{0.995}
\newcommand{\MwPTcemcSmall}{0.38}
\newcommand{\MwLSBblockklBig}{0.001}
\newcommand{\MwPTblockklBig}{0.661}
\newcommand{\MwLSBswtwoBig}{0.18}
\newcommand{\MwPTswtwoBig}{1.41}
\newcommand{\MwLSBruntimeBig}{0.8}
\newcommand{\MwPTruntimeBig}{1.6}
\newcommand{\BgLSBcov}{1.00}
\newcommand{\BgLSBmodes}{6}
\newcommand{\BgPTcov}{0.17}
\newcommand{\BgLSBemc}{0.998}
\newcommand{\BgPTemc}{0.011}
\newcommand{\BgLSBwtwo}{7.82}
\newcommand{\BgBaseWtwo}{14.63}
\newcommand{\BgLSBruntime}{513}}{}
\providecommand{\MogLSBcovK}{n/a}\providecommand{\MogLSBemcK}{n/a}
\providecommand{\MogPTcovK}{n/a}\providecommand{\MogPTemcK}{n/a}
\providecommand{\MogHMCcovK}{n/a}\providecommand{\MogLSBwtwoK}{n/a}
\providecommand{\MogPTwtwoK}{n/a}\providecommand{\MogLSBemcSmall}{n/a}
\providecommand{\MogLSBemcBig}{n/a}\providecommand{\MogLSBruntimeBig}{n/a}
\providecommand{\MogPTruntimeBig}{n/a}\providecommand{\MogLSBpotBig}{n/a}
\providecommand{\MwLSBcemcBig}{n/a}\providecommand{\MwPTcemcBig}{n/a}
\providecommand{\MwLSBcemcSmall}{n/a}\providecommand{\MwPTcemcSmall}{n/a}
\providecommand{\MwLSBblockklBig}{n/a}\providecommand{\MwPTblockklBig}{n/a}
\providecommand{\MwLSBswtwoBig}{n/a}\providecommand{\MwPTswtwoBig}{n/a}
\providecommand{\MwLSBruntimeBig}{n/a}\providecommand{\MwPTruntimeBig}{n/a}
\providecommand{\BgLSBcov}{n/a}\providecommand{\BgLSBmodes}{n/a}
\providecommand{\BgPTcov}{n/a}\providecommand{\BgLSBemc}{n/a}
\providecommand{\BgPTemc}{n/a}\providecommand{\BgLSBwtwo}{n/a}
\providecommand{\BgBaseWtwo}{n/a}\providecommand{\BgLSBruntime}{n/a}

\title{\textbf{L\'evy-Score-Based Monte Carlo Additional Experiments}}
\date{\today}

\begin{document}
\maketitle

\begin{abstract}
We study the problem of drawing samples from a Boltzmann distribution
$\pi(x)\propto\exp(-2V(x)/\sigma^2)$ on $\R^d$ whose mass is concentrated in
several well-separated modes. Gradient-based Markov chain Monte Carlo (MCMC)
mixes well \emph{within} a basin but crosses inter-mode barriers only at a rate
that decays exponentially in barrier height and dimension. \emph{L\'evy-Score-Based
Monte Carlo} (LSB-MC) augments overdamped Langevin dynamics with a
compound-Poisson jump process and a drift correction---the \emph{L\'evy score}---
chosen so that the jump-diffusion retains $\pi$ as its invariant law. We give the
generator of the resulting process and a short proof that the L\'evy-score drift
makes $\pi$ stationary (Proposition~\ref{prop:stationarity}), state the
discretization and quadrature actually used, and define every evaluation metric
precisely. Empirically we benchmark LSB-MC against six samplers, including the two
strongest generic multimodal methods (parallel tempering and Hamiltonian Monte
Carlo), under two scaling axes---mode count $K\in\{10,20,40,80\}$ and dimension
$d\in\{8,16,32,64\}$---and on a Bayesian Gaussian-mixture posterior whose
$K!$ modes arise from an exact label-permutation symmetry. All runs are
single-GPU, seed-replicated with error bars, and charged for compute (wall-clock
and gradient/potential evaluations).
\end{abstract}

% ===========================================================================
\section{Introduction}\label{sec:intro}
Let $\pi(x)\propto e^{-2V(x)/\sigma^2}$ be a target density on $\R^d$. Many tasks
in Bayesian inference and statistical physics reduce to estimating
$\E_\pi[\varphi]$ for test functions $\varphi$, which Monte Carlo addresses by
producing samples $x_1,\dots,x_n$ with empirical law $\hat\pi_n=\frac1n\sum_i
\delta_{x_i}\approx\pi$. The difficulty addressed here is \emph{multimodality}:
$\pi$ has several modes separated by regions of low density. For a diffusion-based
sampler the expected time to move between two modes separated by a barrier of
height $\Delta$ scales (Kramers/Arrhenius) like $\exp(2\Delta/\sigma^2)$, so for
tall barriers or many modes a chain initialised in one basin is, in practice,
absorbed there.

\plainwords{Picture $\pi$ as a landscape whose deep valleys are the likely
outcomes. Sampling drops many walkers and lets them move roughly downhill with
noise, so the final crowd is densest in the deepest valleys. If the valleys are
walled off from each other, walkers that only step locally stay in the valley they
started in and never report the others---so the crowd misrepresents $\pi$.}

\paragraph{Existing remedies and their ceiling.} Tempering-based methods flatten
the barriers in auxiliary high-temperature chains and exchange states downward;
parallel tempering (PT) \cite{geyer1991,earldeem2005} is the strongest
general-purpose instance. Hamiltonian Monte Carlo (HMC) \cite{duane1987,neal2011}
makes long gradient-guided proposals. Both reduce, but do not remove, the
exponential barrier cost: PT needs an increasingly fine temperature ladder as $d$
and the mode count grow, and HMC trajectories still cannot reliably traverse tall
barriers.

\paragraph{The LSB-MC idea.} Rather than crossing barriers by chance, LSB-MC lets
particles take occasional large jumps drawn from a fixed, target-informed
\emph{jump bank}, and adds a drift correction (the L\'evy score) that exactly
offsets the bias such jumps would introduce, preserving $\pi$ as the invariant
law. Section~\ref{sec:method} makes this precise; the correction is the mechanism,
not a tunable add-on (Remark~\ref{rem:nofree}).

\paragraph{Contributions.} (i) We give a self-contained continuous-time
construction and an invariance proof for the L\'evy-score correction
(\S\ref{sec:method}). (ii) We benchmark against the strong baselines PT and HMC,
not only the weak ULA/MALA/FLMC set (\S\ref{sec:samplers}). (iii) We report
interpretable, formally defined metrics---mode coverage, mode-weight KL, effective
mode count (EMC)---alongside $W_2$/MMD (\S\ref{sec:metrics}). (iv) We quantify
scaling in mode count (\S\ref{sec:mog}) and dimension (\S\ref{sec:manywell}), and
study a genuine Bayesian posterior whose modes follow from a symmetry
(\S\ref{sec:bayes}). (v) We account for compute fairly (\S\ref{sec:compute}). The
next page summarises the findings with live numbers.

% ===========================================================================
\section{Results at a glance}\label{sec:glance}
Across all three studies the picture is the same. The local samplers (ULA, MALA, FLMC, ULD, HMC) stay trapped near wherever they were started and never discover the rest of the distribution; parallel tempering (PT), the strongest generic baseline, recovers partial coverage but degrades steadily as the problem grows; and \textbf{LSB-MC reaches essentially the whole distribution}, attaining the best mode coverage and EMC and the smallest mode-weight KL, typically sitting at the finite-sample reference floor.

\textbf{Discovering many modes (MoG).} With $K{=}40$ equally weighted Gaussian modes and all particles started near one of them, LSB-MC covers 1.00 of the modes with EMC 0.94 (1.0 is perfect), while the best baseline (PT) reaches only coverage 0.23 and EMC 0.07; the corresponding Wasserstein distance is 10.00 for LSB-MC versus 33.85 for PT.

\textbf{Scaling to high dimension (ManyWell).} At $d{=}64$ (32 independent double wells), LSB-MC keeps a count-EMC of 0.96 and a per-block marginal KL of 0.001, i.e.\ each well is filled in almost exactly the right proportion, whereas PT and HMC have collapsed to count-EMC $\approx0.00$; the sliced-Wasserstein distance is 0.18 for LSB-MC versus 1.41 for PT.

\textbf{A real Bayesian model (label switching).} The posterior of a 3-component Gaussian mixture has six identical, relabelled copies. Every baseline finds only the one copy it started in (coverage 0.17, 1 of 6 modes, EMC 0.011), whereas LSB-MC populates all six uniformly (coverage 1.00, EMC 0.998) using a jump bank read directly off the label-permutation symmetry.

Compute is accounted for fairly (wall-clock time and gradient evaluations, with PT and HMC charged their extra work and LSB-MC charged its Lévy-score quadrature), so the advantage is not an artefact of spending more compute --- though, as Section~\ref{sec:compute} shows, that quadrature does add a real and problem-dependent per-step cost.

% ===========================================================================
\section{Sampling dynamics and the L\'evy-score correction}\label{sec:method}

\paragraph{Setup and notation.} Throughout, $\pi(x)=Z^{-1}e^{-2V(x)/\sigma^2}$ with
$Z=\int e^{-2V/\sigma^2}\dd x<\infty$, and we fix $\sigma^2=2$ so that
$\pi\propto e^{-V}$, the force is $-\nabla V$, and the diffusive noise has scale
$\sigma\sqrt{\Delta t}$. We write $\nabla\cdot$ for divergence and $\Delta$ for the
Laplacian.

\begin{assumption}\label{ass:reg}
$V\in C^2(\R^d)$ is confining ($e^{-2V/\sigma^2}\in L^1$), and the jump bank is a
finite measure $\nu(\dd r)=\sum_{j=1}^J\lambda_j\,\delta_{r_j}$ on $\R^d$ with total
rate $\lambda=\sum_j\lambda_j<\infty$. These guarantee the generators below act on
$C_c^2$ and the manipulations are justified by dominated convergence.
\end{assumption}

\subsection{Overdamped Langevin diffusion}
The reference dynamics is $\dd X_t=-\nabla V(X_t)\dd t+\sigma\dd B_t$ with
generator and Fokker--Planck adjoint
\[
\mathcal L_0 f=-\nabla V\!\cdot\!\nabla f+\tfrac{\sigma^2}{2}\Delta f,
\qquad
\mathcal L_0^{*}\rho=\nabla\!\cdot\!\big(\nabla V\,\rho\big)+\tfrac{\sigma^2}{2}\Delta\rho .
\]
Since $\nabla\pi=-\tfrac{2}{\sigma^2}\pi\,\nabla V$, i.e.
$\nabla V\,\pi=-\tfrac{\sigma^2}{2}\nabla\pi$, we get
$\mathcal L_0^{*}\pi=\nabla\!\cdot\!(-\tfrac{\sigma^2}{2}\nabla\pi)+\tfrac{\sigma^2}{2}\Delta\pi=0$:
$\pi$ is stationary (indeed reversible). What it lacks is \emph{rate}: barrier
crossings are exponentially rare, which is the failure mode of all the local
samplers in this study.

\subsection{Jump-augmented dynamics}
Add finite-activity jumps with intensity $\nu$ and an auxiliary drift $S$. The
resulting jump-diffusion has generator
\begin{equation}\label{eq:gen}
\mathcal A f(x)=\big(-\nabla V(x)+S(x)\big)\!\cdot\!\nabla f(x)
+\tfrac{\sigma^2}{2}\Delta f(x)
+\int_{\R^d}\!\big[f(x+r)-f(x)\big]\,\nu(\dd r),
\end{equation}
with Fokker--Planck adjoint
\begin{equation}\label{eq:fp}
\mathcal A^{*}\rho(x)=\nabla\!\cdot\!\big((\nabla V-S)\rho\big)
+\tfrac{\sigma^2}{2}\Delta\rho
+\int_{\R^d}\!\big[\rho(x-r)-\rho(x)\big]\,\nu(\dd r).
\end{equation}
For $S\equiv 0$ we get $\mathcal A^{*}\pi=\mathcal L_0^{*}\pi+\int[\pi(\cdot-r)-\pi]\,\nu(\dd r)
=\int[\pi(\cdot-r)-\pi]\,\nu(\dd r)$, which is generically nonzero: bare jumps
destroy stationarity.

\subsection{The L\'evy score and exact invariance}
Define the \emph{L\'evy score}
\begin{equation}\label{eq:score}
S_L(x)\;=\;-\int_0^1\!\!\int_{\R^d} r\,
\exp\!\Big(-\tfrac{2}{\sigma^2}\big(V(x-\theta r)-V(x)\big)\Big)\,\nu(\dd r)\dd\theta
\;=\;-\frac{1}{\pi(x)}\int_{\R^d}\!\!\int_0^1 r\,\pi(x-\theta r)\dd\theta\,\nu(\dd r),
\end{equation}
where the second equality uses $e^{-\frac{2}{\sigma^2}(V(x-\theta r)-V(x))}=\pi(x-\theta r)/\pi(x)$.

\begin{proposition}[Exact stationarity]\label{prop:stationarity}
Under Assumption~\ref{ass:reg}, with $S=S_L$ the jump-diffusion \eqref{eq:gen}
admits $\pi$ as a stationary density: $\mathcal A^{*}\pi=0$.
\end{proposition}
\begin{proof}
Because $\mathcal L_0^{*}\pi=0$, by \eqref{eq:fp} it suffices to show
\begin{equation}\label{eq:key}
\nabla\!\cdot\!\big(S_L\,\pi\big)(x)=\int_{\R^d}\!\big[\pi(x-r)-\pi(x)\big]\,\nu(\dd r).
\end{equation}
From \eqref{eq:score}, the flux is
$S_L(x)\pi(x)=-\int_{\R^d}\!\int_0^1 r\,\pi(x-\theta r)\dd\theta\,\nu(\dd r)$.
Fix $r$. Since $\nabla_x\!\cdot\!\big(r\,\pi(x-\theta r)\big)=r\!\cdot\!\nabla\pi(x-\theta r)$
and $\partial_\theta\,\pi(x-\theta r)=-\,r\!\cdot\!\nabla\pi(x-\theta r)$, we have
$\nabla_x\!\cdot\!\big(r\,\pi(x-\theta r)\big)=-\partial_\theta\,\pi(x-\theta r)$.
Integrating over $\theta\in[0,1]$ telescopes:
\[
\int_0^1\nabla_x\!\cdot\!\big(r\,\pi(x-\theta r)\big)\dd\theta
=-\big[\pi(x-r)-\pi(x)\big]=\pi(x)-\pi(x-r).
\]
Therefore
$\nabla\!\cdot\!(S_L\pi)=-\int_{\R^d}\!\big[\pi(x)-\pi(x-r)\big]\nu(\dd r)
=\int_{\R^d}\!\big[\pi(x-r)-\pi(x)\big]\nu(\dd r)$, which is \eqref{eq:key}.
Substituting into \eqref{eq:fp}, the $-\nabla\!\cdot\!(S_L\pi)$ term cancels the jump
source and $\mathcal A^{*}\pi=0$.
\end{proof}

\begin{remark}[The correction is not an ablation]\label{rem:nofree}
With $S\equiv 0$, $\mathcal A^{*}\pi=\int[\pi(\cdot-r)-\pi]\,\nu(\dd r)\neq 0$ in
general, so the bare jump process is stationary for a \emph{different} law. Hence
``LSB-MC without the L\'evy score'' is not a degraded variant of the same sampler
but a different (incorrect) one; we therefore do not include it as an ablation.
\end{remark}

\begin{remark}[Score interpretation]
The integrand weight $\pi(x-\theta r)/\pi(x)$ is the density ratio along the chord
from $x$ to $x-r$; $S_L$ is the bank-averaged displacement re-weighted by this
ratio. Productive jumps (toward higher $\pi$) and unproductive ones are balanced so
that the net probability flux induced by the jumps is exactly compensated.
\end{remark}

\plainwords{The jumps and the L\'evy-score drift are a matched pair: a new
movement plus the exact bookkeeping that keeps the stationary law equal to $\pi$.
Equation~\eqref{eq:key} is that bookkeeping---the divergence of the corrective flux
equals, pointwise, the net rate at which jumps move probability into each state.}

\subsection{Time discretization and quadrature}
We integrate \eqref{eq:gen} by an explicit Euler--Maruyama step with a
\emph{tamed} drift for numerical stability under stiff $V$ and post-jump
transients. With $\mathrm{tame}(v;\Delta t,c)=\Delta t\,v/(1+(\Delta t/c)\|v\|)$,
one iteration is
\begin{equation}\label{eq:euler}
X^{+}=X+\mathrm{tame}\!\big(-\nabla V(X)+S_L(X);\,\Delta t,\,c\big)+\sigma\sqrt{\Delta t}\,\xi,
\quad \xi\sim\mathcal N(0,I),
\end{equation}
followed by the compound-Poisson jump: draw $N\sim\mathrm{Poisson}(\lambda\Delta t)$
and set $X^{++}=X^{+}+\sum_{m=1}^{N} r_{(m)}$ with each $r_{(m)}\sim\nu/\lambda$
(for the small $\Delta t$ used here $N\in\{0,1\}$ with high probability). The
scheme is weakly consistent of order one with $\mathcal A$; $\pi$ is exactly
stationary only in the continuous-time limit, so a finite $\Delta t$ leaves an
$O(\Delta t)$ bias. Taming (parameter $c$) controls the explosion of the explicit
scheme \cite{hutzenthaler2012}; for stiff targets too large a $c$ overshoots, so
$\Delta t$ and $c$ are chosen jointly (see \S\ref{sec:bayes}).

The L\'evy score \eqref{eq:score} requires the $\theta$-integral, which we evaluate
by $n_\theta$-node Gauss--Legendre quadrature with nodes/weights
$(\theta_q,w_q)$:
\[
S_L(x)\approx-\sum_{q=1}^{n_\theta} w_q\sum_{j=1}^{J}\lambda_j\,r_j\,
\exp\!\Big(-\tfrac{2}{\sigma^2}\big(V(x-\theta_q r_j)-V(x)\big)\Big),
\]
at a cost of $n_\theta J$ potential evaluations per step (charged in
\S\ref{sec:compute}). For ManyWell the integral is separable per coordinate and
precomputed exactly on a fine grid, reducing the per-step cost to $O(1)$ lookups.

\subsection{Jump banks}
The only target-specific ingredient is $\nu$, read off the geometry, never from the
reference: (i) \textbf{ManyWell:} per double-well coordinate $b$,
$\nu=\sum_b \tfrac{\lambda_b}{2}(\delta_{+r e_b}+\delta_{-r e_b})$ with
$r\sim U[3.2,3.8]$ bracketing the well separation; (ii) \textbf{MoG:}
$\nu=\sum_{(i,j)\in\mathrm{kNN}}\lambda\,\delta_{\mu_j-\mu_i}$ over a directed
nearest-neighbour graph of the mixture centres; (iii) \textbf{Bayesian GMM:}
$\nu=\sum_{P}\lambda\,\delta_{P\mu^\star-\mu^\star}$ over the label permutations $P$
applied to a MAP estimate $\mu^\star$.

% ===========================================================================
\section{Samplers and baselines}\label{sec:samplers}
All methods share the target $\pi\propto e^{-V}$ and are run as $S$ independent
particle batches (the seeds). We write $\xi\sim\mathcal N(0,I)$ for fresh noise.

\paragraph{ULA / MALA (local, gradient).}
ULA is \eqref{eq:euler} with $S_L\equiv0$. MALA proposes
$Y=X+\mathrm{tame}(-\nabla V(X);\Delta t,c)+\sigma\sqrt{\Delta t}\,\xi$ and accepts
with probability $\min\{1,\;\pi(Y)q(X\mid Y)/[\pi(X)q(Y\mid X)]\}$, where $q$ is
the Gaussian proposal kernel; this removes the discretization bias at the cost of a
second gradient per step.

\paragraph{ULD (underdamped, BAOAB).}
With momentum $p$, inverse temperature $\beta=2/\sigma^2$ and friction $\gamma$, the
splitting integrator applies, per step, B:$\,p\!\leftarrow\!p-\tfrac{\Delta t}{2}\nabla V(x)$;
A:$\,x\!\leftarrow\!x+\tfrac{\Delta t}{2}p$;
O:$\,p\!\leftarrow\!e^{-\gamma\Delta t}p+\sqrt{(1-e^{-2\gamma\Delta t})/\beta}\,\xi$;
A; B. Its invariant law marginalises to $\pi$ \cite{leimkuhler2013}; momentum
reduces random-walk behaviour but the dynamics remain local.

\paragraph{FLMC (heavy-tailed).}
Langevin driven by isotropic symmetric $\alpha$-stable noise ($\alpha=1.5$):
$X^{+}=X+\mathrm{tame}(-\nabla V(X);\Delta t,c)+c_\alpha\,\Delta t^{1/\alpha}\,L_\alpha$,
$L_\alpha$ an $S\alpha S$ increment. The stationary law solves a fractional
Fokker--Planck equation \cite{simsekli2017}; jumps are long but undirected.

\paragraph{HMC (strong baseline).}
On the augmented Hamiltonian $H(x,p)=V(x)+\tfrac12\|p\|^2$, resample
$p\sim\mathcal N(0,I)$, integrate $L$ leapfrog steps of size $\varepsilon$, and
accept with probability $\min\{1,\,e^{H(x,p)-H(x',p')}\}$. Each proposal costs
$L{+}1$ gradient evaluations.

\paragraph{PT (strong baseline).}
A geometric ladder $1=\beta_1>\beta_2>\dots>\beta_T\ge\beta_{\min}$ with replica $t$
targeting $\pi_t\propto e^{-\beta_t V}$; each replica advances by a (tamed) MALA
move, and every $K_{\mathrm{swap}}$ steps adjacent replicas $(t,t{+}1)$ propose a
swap accepted with probability
$\min\{1,\,\exp[(\beta_t-\beta_{t+1})(V(x_t)-V(x_{t+1}))]\}$. The cold chain
($\beta_1{=}1$) supplies the samples; cost is $T$ kernel evaluations per step.

Step sizes for the Metropolis-based methods (MALA, HMC, PT) are tuned by short
acceptance-rate pilots before the main runs.

% ===========================================================================
\section{Test problems}\label{sec:targets}
Each target exposes $\pi\propto e^{-V}$ and an exact or near-exact reference, so
error is measured against ground truth.

\paragraph{ManyWell (high-dimensional, $d=2n_{\text{blocks}}$).}
$V(x)=\sum_{b=1}^{n_{\text{blocks}}}\big[U(x_{2b-1})+\tfrac12 x_{2b}^2\big]$ with
the double well $U(u)=a u+b u^2+c u^4$, $(a,b,c)=(-0.5,-6,1)$. The two wells per
block are separated by $\approx 3.46$; the bank magnitude $r\sim U[3.2,3.8]$
brackets it. The target has $2^{n_{\text{blocks}}}$ modes; the reference is exact
and i.i.d.\ via the per-coordinate inverse CDF. We sweep
$d\in\{8,16,32,64\}$.

\paragraph{MoG (mode discovery, $d=2$).}
$\pi(x)=\frac1K\sum_{k=1}^{K}\mathcal N(x;\mu_k,I)$ with centres $\mu_k$ fixed
(seed 0) spread over $[-R,R]^2$; all particles start near $\mu_1$, so the task is
to discover and equally weight the remaining $K-1$ modes. The reference is exact
i.i.d. We sweep $K\in\{10,20,40,80\}$.

\paragraph{Bayesian GMM label-switching ($d=6$).}
For data $y_{1:n}$ from a $K{=}3$-component, $2$-D Gaussian mixture with unknown
means $\mu=(\mu_1,\mu_2,\mu_3)$,
\[
V(\mu)=-\sum_{i=1}^{n}\log\frac1K\sum_{k=1}^{K}\mathcal N(y_i;\mu_k,\sigma_y^2 I)
+\sum_{k=1}^{K}\frac{\|\mu_k\|^2}{2\tau^2}.
\]
For any label permutation $P$, $V(P\mu)=V(\mu)$, so the posterior has exactly
$K!{=}6$ identical modes. The bank is the set of differences between permuted MAP
modes; the reference symmetrises a long single-mode chain over all $K!$
permutations.

\plainwords{Calling three clusters $A,B,C$ fits the data exactly as well as
calling them $B,A,C$: the labels are arbitrary. The posterior is therefore six
mirror-image copies, and a correct sampler must visit all six with equal weight.
Reporting only the copy one started in is a valid point estimate but a wrong
posterior.}

% ===========================================================================
\section{Evaluation metrics}\label{sec:metrics}
Let $\hat\pi_n=\frac1n\sum_i\delta_{x_i}$ be the empirical sample law and $\pi$ (or
a large reference sample $Q$) the truth. We use two families: \emph{distributional
distances} (lower is better) and \emph{mode-resolved scores} (coverage/EMC higher
is better).

\paragraph{Maximum mean discrepancy.} For a kernel $k$,
$\MMD^2_k(P,Q)=\E_{x,x'\sim P}k(x,x')+\E_{y,y'\sim Q}k(y,y')-2\,\E_{x\sim P,y\sim Q}k(x,y)$
\cite{gretton2012}. We use a sum of Gaussian kernels
$k(x,y)=\sum_\ell\exp(-\|x-y\|^2/2h_\ell^2)$ with $h_\ell$ multiples of the median
pairwise distance, and report the (biased) $\MMD=\sqrt{\widehat{\MMD^2}}$.

\paragraph{Wasserstein-2.}
$W_2(P,Q)=\big(\inf_{\gamma\in\Pi(P,Q)}\int\|x-y\|^2\dd\gamma\big)^{1/2}$. For MoG and
the Bayesian model we compute the exact empirical $W_2$ by Hungarian assignment on
the $n\times n$ cost matrix \cite{peyre2019}; for high-dimensional ManyWell we use
sliced-$W_2$, $\mathrm{SW}_2^2(P,Q)=\int_{\mathbb S^{d-1}}W_2^2(u_\#P,u_\#Q)\,\dd u$,
estimated over random projections with $1$-D $W_2$ by sorting \cite{bonneel2015}.

\paragraph{Mode-resolved scores.} Assign each sample to its nearest mode (a Voronoi
partition), giving empirical weights $\hat p_m$ against the target weights $p_m^\star$.
Then
\[
\KL(p^\star\|\hat p)=\sum_m p_m^\star\log\frac{p_m^\star}{\max(\hat p_m,1/n)},
\qquad
\mathrm{EMC}=e^{-\KL(p^\star\|\hat p)}\in(0,1],
\]
\[
\text{coverage}=\frac1M\sum_m \mathbf{1}\!\big[\hat p_m>\max(1/n,\,0.001\,p_m^\star)\big].
\]
EMC$=1$ iff the mode weights match; coverage counts the fraction of modes actually
populated.

\paragraph{ManyWell count metrics.} The $2^{n_{\text{blocks}}}$ modes cannot be
enumerated, so we use marginal statistics: the per-block marginal
$\KL$ averaged over blocks, and the law of the number of occupied deep wells
against its $\mathrm{Binomial}(n_{\text{blocks}},p_{\text{deep}})$ target
(count-mode $\KL$, and count-EMC$=e^{-\text{count-}\KL}$), where $p_{\text{deep}}$
is the exact deep-well marginal.

\paragraph{Protocol.} Five seeds, realised as independent particle batches under a
common PRNG stream; we report mean$\pm$SE. Exactly one GPU is used
(\texttt{CUDA\_VISIBLE\_DEVICES} pinned before initialisation); configs and
environment are saved per run.

% ===========================================================================
\section{Experiment 1: scaling in the number of modes (MoG)}\label{sec:mog}
\paragraph{Setup.} $K$ equally weighted Gaussian modes, all particles started near
$\mu_1$; a correct sampler reaches coverage and EMC near $1$ with small $W_2$/MMD.
Because the start is concentrated, this isolates inter-mode \emph{transport}.

\paragraph{Results.} Table~\ref{tab:mog_emc_scaling} reports EMC across the sweep
and Table~\ref{tab:mog_K40} the full metric set at $K{=}40$;
Figure~\ref{fig:mog_conv} tracks mode coverage over sampling time. At $K{=}40$,
LSB-MC attains coverage \MogLSBcovK{} and EMC \MogLSBemcK, versus PT's coverage
\MogPTcovK{} / EMC \MogPTemcK{} and HMC's coverage \MogHMCcovK; $W_2$ is
\MogLSBwtwoK{} (LSB-MC) versus \MogPTwtwoK{} (PT). LSB-MC's EMC stays high across
the sweep (\MogLSBemcSmall{} at $K{=}10$, easing to \MogLSBemcBig{} at $K{=}80$),
while the local methods never leave $\mu_1$.

\figbox{mog_convergence}{MoG mode coverage over sampling time (mean over
seeds).}{fig:mog_conv}
\inputtab{mog_emc_scaling}
\inputtab{mog_table_K40}

\paragraph{Discussion.} The gap is qualitative, not incremental. $W_2$ is
non-monotone in $K$ for LSB-MC---an artefact of comparing point clouds under large
separation---so coverage, EMC and MMD are the more reliable measures here; they
agree that LSB-MC recovers the mixture while the baselines do not.

% ===========================================================================
\section{Experiment 2: scaling in dimension (ManyWell)}\label{sec:manywell}
\paragraph{Setup.} $n_{\text{blocks}}$ independent double wells give
$2^{n_{\text{blocks}}}$ modes; we sweep $d\in\{8,16,32,64\}$. Independence makes the
exact answer simple to state (each block deep a fixed fraction of the time), which
lets us measure error cleanly despite the combinatorial mode count.

\paragraph{Results.} Table~\ref{tab:manywell_countemc_scaling} reports count-EMC
across dimension and Table~\ref{tab:manywell_d64} the full metric set at $d{=}64$;
Figure~\ref{fig:mw_conv} tracks count-EMC over sampling time. At $d{=}64$, LSB-MC
has count-EMC \MwLSBcemcBig{} and per-block marginal $\KL$ \MwLSBblockklBig, versus
PT's count-EMC \MwPTcemcBig{} and block $\KL$ \MwPTblockklBig; the sliced-$W_2$ is
\MwLSBswtwoBig{} (LSB-MC) versus \MwPTswtwoBig{} (PT). PT starts respectably
(count-EMC \MwPTcemcSmall{} at $d{=}8$) but decays toward zero, since getting all
blocks right simultaneously is exponentially unlikely without directed jumps.

\figbox{manywell_convergence}{ManyWell count-EMC over sampling time (mean over
seeds).}{fig:mw_conv}
\inputtab{manywell_countemc_scaling}
\inputtab{manywell_table_d64}

\paragraph{Discussion.} Because the bank acts per block and brackets the
separation, every barrier is crossed productively and LSB-MC pays essentially no
penalty for added dimension, while undirected exploration decays as the modes
multiply.

% ===========================================================================
\section{Experiment 3: a Bayesian posterior with a symmetry (label switching)}\label{sec:bayes}
\paragraph{Setup.} Inferring three exchangeable cluster means yields a posterior
with six identical, label-permuted modes (\S\ref{sec:targets}); a correct sampler
must weight all six equally.

\paragraph{Results.} Table~\ref{tab:bayes_gmm} reports the full metric set;
Figure~\ref{fig:bg_conv} tracks mode coverage over sampling time. Every
baseline---including PT and HMC---stays in its initial mode (coverage \BgPTcov,
$1$ of $6$, EMC \BgPTemc), whereas LSB-MC populates all six uniformly (coverage
\BgLSBcov{} $=\BgLSBmodes/6$, EMC \BgLSBemc) and is closest in earth-mover distance
($W_2=\BgLSBwtwo$ versus $\approx\BgBaseWtwo$ for the trapped baselines).

\figbox{bayes_gmm_convergence}{Bayesian GMM mode coverage over sampling time (mean
over seeds).}{fig:bg_conv}
\inputtab{bayes_gmm_table}

\paragraph{Discussion.} On a genuine posterior, exploring the symmetry is not
optional: reporting one mode discards five-sixths of the answer. LSB-MC recovers
the full posterior from a bank that needed no tuning beyond writing down the
symmetry. As Remark~\ref{rem:nofree} and \S\ref{sec:method} predict, discretization
matters: an over-aggressive jump schedule strands particles between modes, and the
configuration used here (finer $\Delta t$, more steps, moderate $\lambda$) clears
that transit mass, returning the posterior-energy distribution to the reference's
low-energy band.

% ===========================================================================
\section{Compute-fair comparison}\label{sec:compute}
A faster-mixing sampler is only interesting if the mixing is not bought with hidden
compute. Table~\ref{tab:compute} reports, at the largest setting of each study,
wall-clock time, gradient evaluations, and L\'evy-score potential evaluations, with
HMC charged $L{+}1$ gradients per proposal, PT charged $T$ kernels per step, and
LSB-MC charged its quadrature potential evaluations. Per-setting time and gradient
counts also appear in the result tables of \S\ref{sec:mog}--\S\ref{sec:bayes}.

\inputtab{compute_table}

\paragraph{Reading.} LSB-MC's cost is dominated by the quadrature and is
problem-dependent. On ManyWell the correction is a $1$-D grid lookup, so LSB-MC is
the \emph{cheapest} method ($\approx\MwLSBruntimeBig$\,s at $d{=}64$ versus
$\MwPTruntimeBig$\,s for PT) while being the most accurate---it wins on both axes.
On MoG the correction integrates over a large nearest-neighbour bank, which is
genuinely expensive at $K{=}80$ ($\approx\MogLSBruntimeBig$\,s and
$\approx\MogLSBpotBig$\,M potential evaluations, versus $\MogPTruntimeBig$\,s for
PT): there the quality win carries a real per-step cost. The Bayesian model's
$\BgLSBruntime$\,s reflects the deliberately fine step size, not an intrinsic floor.
The advantage is therefore not an artefact of extra compute, but on dense banks the
quadrature is the price of admission.

% ===========================================================================
\section{Discussion}\label{sec:discussion}
The three experiments vary, respectively, the mode count, the dimension, and the
\emph{origin} of the modes (a real symmetry). Across all three the ordering is
stable: local methods stay put; PT helps but degrades as the problem grows; LSB-MC
tracks the reference. The mechanism is the one made precise in \S\ref{sec:method}:
undirected crossings cost $\exp(2\Delta/\sigma^2)$, which compounds with mode count
and dimension, whereas LSB-MC replaces them with directed jumps whose bias is
cancelled exactly (Proposition~\ref{prop:stationarity}), decoupling mixing from
barrier height.

% ===========================================================================
\section{Limitations}\label{sec:limitations}
\begin{itemize}
\item \textbf{The bank needs structure.} $\nu$ is informed by the target geometry;
this suits structured problems but is not black-box, whereas PT/HMC need no such
input. Constructing good banks when the geometry is unknown is open.
\item \textbf{No universal dominance.} Where a baseline is competitive the tables
show it (e.g.\ FLMC on low-$d$ ManyWell).
\item \textbf{Dense banks cost compute} (\S\ref{sec:compute}).
\item \textbf{Continuous-time theory, discrete implementation.}
Proposition~\ref{prop:stationarity} is exact for $\mathcal A$; the Euler/quadrature
scheme leaves an $O(\Delta t)$ bias and can strand particles if mis-tuned.
\item \textbf{High-$d$ $W_2$ is a proxy} (sliced-$W_2$).
\end{itemize}

% ===========================================================================
\section{Conclusion}\label{sec:conclusion}
Multimodal sampling is hard because modes are walled off and undirected steps pay
an exponential price to cross. LSB-MC instead jumps along target-informed
directions and corrects the induced bias exactly, so that $\pi$ remains the
invariant law (Proposition~\ref{prop:stationarity}). Empirically, under strong
baselines, interpretable metrics, two scaling axes, and a genuine Bayesian
posterior, LSB-MC reaches and correctly weights essentially the whole distribution
where the baselines stall---at a compute cost that is honestly reported and, on the
high-dimensional benchmark, lower than the baselines'. The one requirement is a
structure-informed jump bank, which is exactly what the structured, hardest-to-mix
problems provide.

% ===========================================================================
\section{Recommendations for the main paper}\label{sec:recs}
% Recommended manuscript changes (static guidance; numbers in the tables).
\paragraph{Main paper.} Replace the single-target showcase with two pieces of
evidence that referees ask for: (i) a \emph{scaling} result and (ii) a
\emph{realistic} benchmark.
\begin{itemize}
\item Promote the \textbf{MoG mode-count scaling} figure (Fig.~\ref{fig:mog_scaling})
and the \textbf{ManyWell dimension-scaling} figure (Fig.~\ref{fig:mw_scaling}) to
the main paper: they show how coverage / EMC degrade with $K$ and $d$ for the
local baselines while LSB-MC stays close to the reference.
\item Add the \textbf{Bayesian GMM label-switching} experiment
(Fig.~\ref{fig:bg_modeweights}, Table~\ref{tab:bayes_gmm}) as the headline
``real model'' result: the jump bank is derived from a genuine symmetry of the
posterior rather than hand-tuned, which strengthens the methodological story.
\item Include the strongest baselines (PT and HMC), not just ULA/MALA/FLMC. The
honest comparison against PT is what makes the multimodal claim credible.
\item Show one \textbf{compute-fairness} panel (quality vs.\ wall-clock /
gradient evaluations) so the speed-up is not attributed to extra compute.
\end{itemize}

\paragraph{Appendix.} Per-dimension and per-$K$ tables, acceptance/swap
diagnostics, the exact configurations, and the reference-construction details for
the Bayesian posterior.

\paragraph{What the old section should become.} The original double-well /
single-MoG / single-ManyWell demonstrations are good \emph{illustrations} and
should move to the appendix or a ``setup'' figure; the quantitative claims should
rest on the scaling studies and the Bayesian benchmark.


% ===========================================================================
\begin{thebibliography}{9}
\bibitem{geyer1991} C.~J.~Geyer. Markov chain Monte Carlo maximum likelihood.
\emph{Computing Science and Statistics: Proc.\ 23rd Symposium on the Interface}, 1991.
\bibitem{earldeem2005} D.~J.~Earl and M.~W.~Deem. Parallel tempering: theory,
applications, and new perspectives. \emph{Phys.\ Chem.\ Chem.\ Phys.}, 7:3910--3916, 2005.
\bibitem{duane1987} S.~Duane, A.~D.~Kennedy, B.~J.~Pendleton, D.~Roweth. Hybrid
Monte Carlo. \emph{Physics Letters B}, 195(2):216--222, 1987.
\bibitem{neal2011} R.~M.~Neal. MCMC using Hamiltonian dynamics. \emph{Handbook of
Markov Chain Monte Carlo}, CRC Press, 2011.
\bibitem{leimkuhler2013} B.~Leimkuhler and C.~Matthews. Rational construction of
stochastic numerical methods for molecular sampling. \emph{Appl.\ Math.\ Res.\
Express}, 2013(1):34--56, 2013.
\bibitem{hutzenthaler2012} M.~Hutzenthaler, A.~Jentzen, P.~E.~Kloeden. Strong
convergence of an explicit numerical method for SDEs with non-globally Lipschitz
continuous coefficients. \emph{Ann.\ Appl.\ Probab.}, 22(4):1611--1641, 2012.
\bibitem{simsekli2017} U.~\c{S}im\c{s}ekli. Fractional Langevin Monte Carlo:
exploring L\'evy driven stochastic differential equations for MCMC.
\emph{ICML}, 2017.
\bibitem{gretton2012} A.~Gretton, K.~Borgwardt, M.~Rasch, B.~Sch\"olkopf, A.~Smola.
A kernel two-sample test. \emph{JMLR}, 13:723--773, 2012.
\bibitem{bonneel2015} N.~Bonneel, J.~Rabin, G.~Peyr\'e, H.~Pfister. Sliced and
Radon Wasserstein barycenters of measures. \emph{J.\ Math.\ Imaging Vis.},
51:22--45, 2015.
\bibitem{peyre2019} G.~Peyr\'e and M.~Cuturi. Computational optimal transport.
\emph{Found.\ Trends Mach.\ Learn.}, 11(5--6):355--607, 2019.
\bibitem{applebaum2009} D.~Applebaum. \emph{L\'evy Processes and Stochastic
Calculus}. Cambridge University Press, 2nd edition, 2009.
\end{thebibliography}

% ===========================================================================
\appendix
\section{Configurations and hyperparameters}
\small
\paragraph{manywell scaling}
\begin{itemize}
\item scaling: \texttt{\detokenize{n_blocks}} $\in$ [4, 8, 16, 32]
\item methods: ULA, MALA, FLMC, LSBMC, ULD, HMC, PT
\item particles 2048, steps 2000, dt 0.005, seeds [0, 1, 2, 3, 4]
\item method configs: \texttt{\detokenize{FLMC: {'alpha': 1.5, 'dt': 0.005}; HMC: {'dt': 0.05, 'n_leapfrog': 10}; LSBMC: {'dt': 0.005}; MALA: {'dt': 0.005}; PT: {'beta_min': 0.05, 'dt': 0.005, 'n_temps': 6, 'swap_interval': 5}; ULA: {'dt': 0.005}; ULD: {'dt': 0.01, 'friction': 2.0}}}
\item target: \texttt{\detokenize{{'jump_r_max': 3.8, 'jump_r_min': 3.2, 'lambda_block': 0.1}}}
\item output: \texttt{\detokenize{results/iclr_sampling/20260628_032655_manywell_scaling}}
\end{itemize}
\paragraph{mog scaling}
\begin{itemize}
\item scaling: \texttt{\detokenize{n_modes}} $\in$ [10, 20, 40, 80]
\item methods: ULA, MALA, FLMC, LSBMC, ULD, HMC, PT
\item particles 1500, steps 1500, dt 0.01, seeds [0, 1, 2, 3, 4]
\item method configs: \texttt{\detokenize{FLMC: {'alpha': 1.5, 'dt': 0.01}; HMC: {'dt': 0.1, 'n_leapfrog': 10}; LSBMC: {'dt': 0.01, 'tame_cap': 10.0}; MALA: {'dt': 0.01}; PT: {'beta_min': 0.02, 'dt': 0.01, 'n_temps': 8, 'swap_interval': 5}; ULA: {'dt': 0.01}; ULD: {'dt': 0.02, 'friction': 2.0}}}
\item target: \texttt{\detokenize{{'center_range': 40.0, 'center_seed': 0, 'comp_sigma': 1.0, 'jump_chunk': 128, 'k_neighbors': 8, 'lam_total': 1.0, 'n_theta': 8}}}
\item output: \texttt{\detokenize{results/iclr_sampling/20260628_032158_mog_scaling}}
\end{itemize}
\paragraph{bayes gmm label switching}
\begin{itemize}
\item scaling: \texttt{\detokenize{K}} $\in$ [3]
\item methods: ULA, MALA, FLMC, LSBMC, ULD, HMC, PT
\item particles 2000, steps 2000, dt 0.01, seeds [0, 1, 2, 3, 4]
\item method configs: \texttt{\detokenize{FLMC: {'alpha': 1.5, 'dt': 0.01}; HMC: {'dt': 0.05, 'n_leapfrog': 15}; LSBMC: {'dt': 0.004, 'metric_every': 250, 'n_steps': 5000}; MALA: {'dt': 0.01}; PT: {'beta_min': 0.02, 'dt': 0.01, 'n_temps': 8, 'swap_interval': 5}; ULA: {'dt': 0.01}; ULD: {'dt': 0.02, 'friction': 2.0}}}
\item target: \texttt{\detokenize{{'data_seed': 0, 'jump_chunk': 32, 'lam_total': 0.4, 'n_data': 300, 'n_theta': 8, 'p': 2, 'sigma_y': 1.0, 'tau': 10.0, 'true_radius': 6.0}}}
\item output: \texttt{\detokenize{results/iclr_sampling/20260628_031115_bayes_gmm_label_switching}}
\end{itemize}

\section{Additional per-setting tables}
\inputtab{mog_table_K10}\inputtab{mog_table_K20}\inputtab{mog_table_K80}
\inputtab{manywell_table_d8}\inputtab{manywell_table_d16}\inputtab{manywell_table_d32}

\section{Acceptance / swap diagnostics}
\inputtab{acceptance_table}

\end{document}
